\section{Question 1: étude de brevetabilité}

\subsection{Droit opposable}

La demande de brevet EP a été déposée le 26/10/2020.
Donc, le droit applicable est la \textbf{CBE 2000}.
\begin{itemize}
  \item \epcart{99}(1) CBE dispose que le délai d'opposition est de 9 mois après le 26/02/2025, 
  date de la publication de la mention de la mention de délivrance
  du brevet EP au Bulletin Européen des Brevets;
  
  \item \epcrule{131}(4) CBE donne que le délai d'opposition expire le mercredi 26/11/2025;

  \item \epcrule{134}(1) CBE ne s'applique pas, le 26/11/2025 étant un mercredi.
\end{itemize}
Donc, le 8/03/2025 est dans le délai d'opposition.\\

\textbf{Conclusion: les motifs de \epcart{100} CBE sont susceptibles d'être invoqués à l'encontre du brevet EP.}\\

Le motif de \epcart{100}c) CBE ne s'applique pas.
En effet, le texte du brevet EP délivré correspond à la demande telle que déposée (énoncé paragraphe 4 page 1).
Le motif de \epcart{100}b) CBE n'est pas valable a priori.
Nous allons examiner le motif de \epcart{100}a) CBE.

\subsection{Divulgation opposable}

La demande de brevet EP a été déposée le 26/10/2020 sans priorité.
Donc, la date effective retenue est le 26/10/2020 pour toutes les revendications. \\

\begin{itemize}
    \item D1 est un brevet FR publié le 25/11/1931;
    \item D2 est une demande de brevet US publiée le 28/02/2019;
    \item D3 est une demande de brevet EP publiée le 23/06/2010.
\end{itemize}
Donc, D1, D2 et D3 sont publiées avant la date effective.\\

\textbf{Conclusion: D1, D2 et D3 sont des états de la technique \epcart{54}(2) CBE de la demande de brevet EP.}\\

NL308' est une demande néerlandaise déposée le 20/04/2020 avant la date effective et publiée après
la date effective. Mais, NL308' n'est pas une demande européenne.
Donc, NL308' n'est pas un état de la technique \epcart{54}(2) CBE ou \epcart{54}(3) CBE de la demande de brevet EP.
En revanche, NL308' est un droit antérieur national
pour la partie néerlandaise du brevet EP.

\newpage

\subsection{Nouveauté}

\subsubsection{Nouveauté par rapport à D1}

D1 divulgue un instrument à commande manuelle
("manœuvrable à la main" p.1 col.2 l.35) comportant :
\begin{itemize}
  \item un piston ("rondelle de butée 5" p.1 col.2 l.36);
  \item un axe fileté ("tige filetée 2" p.1 col.2 l.33)
  portant un filetage extérieur (figures 1, 4, 5 et 6);
  \item un écrou d'actionnement manuel ("écrou 4 manœuvrable à la main"
  p.1 col.2, l.35) pourvu d'un trou traversant ("traversé librement
  par la tige filetée 2" p.1 col.2 l.43).\\
\end{itemize}

D1 ne divulgue pas:
\begin{itemize}
  \item une seringue;
  \item un siège adapté pour recevoir un axe
  de commande du coulissement du piston à l'intérieur du corps de la seringue
  \item un moyen d'accouplement au siège du piston;
  \item une cinématique telle que lorsque le moyen d'accouplement est couplé au siège du piston, la 
  rotation de l'écrou d'actionnement par rapport au corps de la seringue provoque 
  le coulissement axial de l'axe fileté au travers de son trou traversant et le déplacement 
  du piston à l'intérieur du corps de la seringue.\\
\end{itemize}

D1 divulgue un procédé de mise en œuvre d'un instrument à commande manuelle
("manœuvrable à main" p.1 col.2 l.35).\\

D1 ne divulgue pas:
\begin{itemize}
  \item une seringue comprenant un piston pourvu d'un siège;
  \item une étape de couplage de l'instrument à la seringue, définie selon la revendication 3;
  \item un déplacement de l'écrou d'actionnement, défini selon la revendication 3;
  \item une rotation de l'écrou d'actionnement, définie selon la revendication 3.\\
\end{itemize}

Par ailleurs, la revendication 2 dépend de la revendication 1.\\

\textbf{Conclusion: les revendications 1, 2 et 3 sont nouvelles par rapport à D1.}\newpage

\subsubsection{Nouveauté par rapport à D2}

D2 divulgue un instrument à commande manuelle pour seringue
("seringe assembly 10" p.1 l.29) comportant:
\begin{itemize}
  \item un piston ("plunger 24" p.1 l.37) pourvu d'un siège ("threaded bore 28" p.1 l.42)
  adapté pour recevoir un axe de commande ("plunger shaft 32" p.1 l.45) du coulissement
  du piston à l'intérieur du corps de la seringue ("extend rearwardly away from
  the plunger centrally of the interior of the barrel" p.1 l.48),
\end{itemize}
 l'instrument à commande manuelle comportant:
\begin{itemize}
  \item un moyen d'accouplement au siège du piston ("a threaded extension 34" p.1 l.46);
  \item un axe fileté portant un filetage extérieur ("a spiral rib 38 is formed
  on the external surface of the plunger shaft 32" p.2 l.2-3);
  \item un écrou d'actionnement manuel pourvu d'un trou traversant portant un filetage intérieur
  ayant le même pas que le filetage de l'axe fileté et au travers duquel l'axe fileté
  passe et vient en prise ("a threaded bore 46 extends centrally of the nut with
  an internal spiral groove to threadly engage the spiral rib 38" p.1 l.11-12).\\
\end{itemize}

D2 ne divulgue pas:
\begin{itemize}
  \item un axe fileté en extrémité duquel le moyen d'accouplement est \underline{solidement fixé};
  \item une cinématique telle que définie à la revendication 1.\\
\end{itemize}

En effet, \epcart{69}(1) CBE dispose que les revendications sont interprétées en fonction de la description. Ainsi,
le terme "solidement fixé" s'interprète à la lumière de la description:
"\textit{De préférence, les moyens d'accouplement 13 de l'instrument 1 au siège 7 du piston 5 comprennent
une tête filetée 13a pour cette fixation qui soit suffisamment solide compte tenu, d'une part,
de l'effort de traction exercé par l'axe 9 qui contrôle le déplacement du piston 5 et, d'autre part,
de la résistance exercée à l'encontre du coulissement de ce piston 5.}" (par.[0011]). Or,
dans D2, c'est précisément l'inverse qui est décrit: "\textit{as the plunger shaft 32 is advanced through 
the nut and caused to rotate in a clockwise direction will impart rotation to the extension 34 
in a direction causing   it to be unthreaded or released from the plunger 24, since again the
plunger 24 is restrained against rotation by virtue of its frictional engagement with
the inner wall surface of the barrel 12.}" (p.2 l.30-32).
Donc, D2 ne divulgue pas un axe fileté
en extrémité duquel le moyen d'accouplement est \underline{solidement fixé}.\\

% De même, par \epcart{69}(1) CBE, le terme "actionnement manuel" s'interprète à la lumière de la description par 
% "\textit{mouvement rotatif de l'écrou d'actionnement 15}" (par.[0013]).\\
% En revanche, dans D2, l'écrou présente une position angulaire fixe:
% "\textit{the nut 40 is inserted into the flanged end portion 18 of the barrel and permanently
% attached as described so that any further movement of the plunger shaft 32 is solely by rotation
% through the nut 40.}" (p.2 l.18-20).\\
% Donc, D2 ne divulgue pas un écrou \underline{d'actionnement manuel}.\\

D2 divulgue un procédé de mise en œuvre d'un instrument à commande manuelle avec une seringue
comprenant un piston pourvu d'un siège, le procédé comprenant:
\begin{itemize}
  \item une étape de couplage de l'instrument à la seringue par le couplage
  du moyen d'accouplement au siège du piston ("the plunger 24 is threaded onto the end
  of the plunger shaft 32" p.2 l.15);
  \item un déplacement de l'écrou d'actionnement en butée contre le bord de l'ouverture du corps
  de la seringue au niveau de laquelle le piston entre dans le corps de la seringue
  ("the nut 40 is inserted into the flanged end portion 18 of the barrel" p.2 l.18-19).\\
\end{itemize}

D2 ne divulgue pas:
\begin{itemize}
  \item une rotation de l'écrou d'actionnement, définie selon la revendication 3.\\
\end{itemize}

\textbf{Conclusion: les revendications 1, 2 et 3 sont nouvelles par rapport à D2.}\newpage

\subsubsection{Nouveauté par rapport à D3}

D3 divulgue un instrument à commande manuelle pour seringue ("hypodermic syringe
and dose-dispenser device" p.1 l.34-35) comportant:
\begin{itemize}
  \item un piston pourvu d'un siège adapté pour recevoir un axe de commande du coulissement
  du piston à l'intérieur du corps de la seringue ("positioned within cartridge 2, for
  axial reciprocation therein, is a plunger 15 which, on its surface
  facing toward the distal open end of cartridge 2, has a screw-threaded extension 16" 
  p.1 l.43-45),
\end{itemize}
l'instrument à commande manuelle comportant:
\begin{itemize}
  \item un moyen d'accouplement au siège du piston ("a plunger screw-mating internal bore 17 by which rod
  10 and plunger 15 are connectable" p.1 l.45-46);
  \item un axe fileté portant un filetage extérieur ("the combined needle cover and plunger rod 10 is
  further provided over the major portion of its surface with an external thread 18)
   et en extrémité duquel le moyen d'accouplement est solidement fixé ("combined needle cover and plunger
  push rod 10 has at its proximal end a plunger screw-mating internal bore 17" p.1 l.45-46);
  \item un écrou d'actionnement manuel pourvu d'un trou traversant portant un filetage
  intérieur ayant le même pas que le filetage de l'axe fileté ("a sleeve 19 is rotatably mounted
  by means of its internally mating thread 20" p.1 l.2-3).\\
\end{itemize}

D3 ne divulgue pas:
\begin{itemize}
  \item une cinématique de sorte que, lorsque le moyen d'accouplement est couplé au siège du piston,
  la rotation de l'écrou d'actionnement par rapport au corps de la seringue provoque le coulissement axial
  de l'axe fileté au travers de son trou traversant.\\
\end{itemize}

D3 divulgue un procédé de mise en œuvre d'un instrument à commande manuelle avec une seringue
comprenant un piston pourvu d'un siège, le procédé comprend:
\begin{itemize}
  \item une étape de couplage de l'instrument à la seringue par le couplage du moyen d'accouplement
  au siège du piston;
  \item un déplacement de l'écrou d'actionnement en butée contre le bord de l'ouverture du corps
  de la seringue au niveau de laquelle le piston entre dans le corps de la seringue
  ("sleeve 19 is rotated on rod 10 to cause the former to be positioned with its frusto-conical
  surface 22 abutting within the flared inner distal end 3 of cartridge 2." p.2 l.35-36).\\
\end{itemize}

D3 ne divulgue pas:
\begin{itemize}
  \item une rotation de l'écrou d'actionnement en butée par rapport au corps de la seringue,
  provoquant ainsi le coulissement axial de l'axe fileté à l'intérieur du trou de l'écrou
  d'actionnement et une traction du piston vers l'ouverture d'extrémité du corps de la seringue
  pour réaliser une aspiration au travers de la pointe de la seringue.\\
\end{itemize}

\textbf{Conclusion: les revendications 1, 2 et 3 sont nouvelles par rapport à D3.}\\

\newpage

\subsection{Activité inventive}

\subsubsection{Etat de la technique le plus proche}

L'invention fait partie du domaine technique des seringues. L'invention cherche à faciliter
le remplissage manuel des seringues.\\

D1 fait partie du domaine technique des tire-bouchons. D1 cherche à enlever facilement le bouchon
de la mèche (p.1 col.1 l.1-4).\\

D2 fait partie du domaine technique des seringues. D2 cherche à éviter le remplissage d'une seringue
après utilisation (p.1 l.12-13).\\

D3 fait partie du domaine technique des seringues. D3 cherche un moyen de protection de l'aiguille 
d'une seringue et d'administration de doses précises (p.1 l.14-17).\\

D3 fait partie du même domaine technique que l'invention et cherche à résoudre le même problème technique
que l'invention.\\

\textbf{Conclusion: D3 est l'état de la technique le plus proche.}\\

\subsubsection{Problème technique objectif}

Une caractéristique distinctive de la revendication 1 par rapport à D3 est une cinématique de sorte que, 
lorsque le moyen d'accouplement est couplé au siège du piston,
  la rotation de l'écrou d'actionnement par rapport au corps de la seringue provoque le coulissement axial
  de l'axe fileté au travers de son trou traversant.\\

L'effet technique de cette caractéristique distinctive est d’obtenir 
"un effort de traction sur le piston" (par.[0013]).\\

Le problème technique objectif est donc: comment obtenir un effort de traction sur le piston d'une seringue?\\

Le même raisonnement s'applique mutatis mutandis au procédé défini à la revendication 3. 

\subsubsection{Résolution}

La personne du métier est spécialiste en seringues.\\

D3 résout déjà le problème technique objectif: "plunger 15 may be reciprocated within cartridge 2 by
axial manipulation push rod 10 to permit aspiration or reconstitution of the liquid L therein, or,
in certain instances, to take up liquid initially for charging the cartridge." (p.2 l.25-27).\\

Ainsi, la personne du métier n'aurait pas été incitée à modifier D3 pour parvenir à la solution
définie à la revendication 1.\\

D3 présente un enseignement contraire à l'abuttement de l'écrou d'actionnement par rapport au corps
de la seringue défini à la revendication 3 pour le remplissage de la seringue:
"sleeve 19 may have to be rotated on rod 10 to position the former
on the latter toward the distal end so as not to restrict the required inward stroke of
rod 10 by premature abutment of sleeve 19 against the flared end 3 of cartridge 2" (p.2 l.28-30).\\

Même si pour une raison inconnue, la personne du métier place l'écrou de D3 en butée par rapport au corps
de seringue, l'écrou de D3 n'est pas en mesure d'effectuer un mouvement de rotation:
"due to the frusto-conical configuration of surface 22 of sleeve 19, the sleeve 19 is locked against
rotation within said flared end when abutted or inserted within." (p.2 l.36-38). Donc D3 présente
une incompatibilité technique par rapport à la solution de la revendication 3.\\

Ainsi, la personne du métier, face au problème technique objectif,
n'aurait pas été incitée et n'aurait pas pu parvenir à modifier D3 pour 
arriver à la solution définie à la revendication 3.\\

\textbf{Conclusion: les revendications 1, 2 et 3 impliquent une activité inventive par rapport à D3}\\

\subsection{Conclusion}

\textbf{Aucun motif d'opposition \epcart{100} CBE ne s'applique aux revendications 1 à 3. Les chances de succès
à une opposition au brevet EP01 sont élevées.}